% DROP-IN replacement for the appendix TOST definition.
% The original defined theta = |mu2 - mu1| * eps -- a margin derived
% from the observed difference -- and specified independent samples.
% Generated by recompute/make_report.py.
\subsection*{TOST test}

A two one-sided t-test (TOST) assesses whether two systems are equivalent to within a margin fixed in advance. Because the same queries are run through both systems, the test is applied to the \emph{paired} per-query differences $d_q = m(\text{candidate}, q) - m(\text{baseline}, q)$ rather than to the two score sets as independent samples; pairing removes the between-query variance that dominates retrieval effectiveness.

We fix an equivalence margin $\theta = \varepsilon\,\mu_{\text{baseline}}$, a proportion $\varepsilon$ of the baseline's mean effectiveness, with $\varepsilon = 0.01$ and $\alpha = 0.05$. Note that $\theta$ is defined relative to the baseline mean and is chosen before the comparison is made; it is not a function of the observed difference. We then test two one-sided hypotheses, $H_{01}: \mathbb{E}[d] \leq -\theta$ and $H_{02}: \mathbb{E}[d] \geq \theta$, each at level $\alpha$ using a paired $t$-test. Equivalence is concluded only if \emph{both} are rejected, which is identical to requiring that the $90\%$ confidence interval of $\mathbb{E}[d]$ lie entirely within $(-\theta, \theta)$. Rejecting both at level $\alpha$ gives an overall type-I error rate of $\alpha$, not $2\alpha$.

Failing to conclude equivalence is not evidence of a difference, and a non-significant difference is not evidence of equivalence; the two questions are separate, so we also report a two-sided paired $t$-test. For each comparison we report $\varepsilon^{*} = \max(|\text{CI}_{\text{low}}|, |\text{CI}_{\text{high}}|) / \mu_{\text{baseline}}$, the smallest margin the data support, so that a reader may apply any margin they consider material.
